\section{Introduction}\label{a:sec1}

Let $\geh$ be a f\/inite-dimensional simple Lie algebra and
$U_q(\geh)$ be the Drinfeld--Jimbo quantized enveloping algebra.
$U_q(\geh)$ has the subalgebra $U^+_q(\geh)$ generated by the
Chevalley generators $e_1,\ldots, e_n$   $(n= \mathrm{rank}\, \geh)$
corresponding to the simple roots.
Denote by $W=\langle s_1,\ldots, s_n\rangle$ the
Weyl group of $\geh$ generated by the simple ref\/lections $s_1,\ldots, s_n$.
It is well known (see for example~\cite{L2}) that
for each reduced expression $w_0 = s_{i_1}\cdots s_{i_l}$
of the longest element $w_0 \in W$,
one can associate the Poincar\'e--Birkhof\/f--Witt (PBW) basis
of $U^+_q(\geh)$ having the form
\begin{gather*}
E_{\bf i}^A=e_{\beta_1}^{(a_1)}e_{\beta_2}^{(a_2)}\cdots e_{\beta_l}^{(a_l)},
\qquad
A=(a_1,\ldots, a_l) \in (\Z_{\ge 0})^l,
\end{gather*}
where $e^{(a_i)}_{\beta_i}$'s are the divided powers of the positive root
vectors determined by the choice ${\bf i}= (i_1,\ldots, i_l)$.
See Section~\ref{a:ss:pbw}.
Let $\big\{E^A_{\bf j}\mid A \in (\Z_{\ge 0})^l\big\}$ with
${\bf j} = (j_1,\ldots, j_l)$ be another
PBW basis associated with a yet dif\/ferent reduced expression
$w_0 = s_{j_1}\cdots s_{j_l}$.
Following Lusztig~\cite{L}, one expands
a basis in terms of another as
\begin{gather*}
E^A_{\bf i} = \sum_{B \in (\Z_{\ge 0})^l}
\gamma^A_B E^B_{\bf j}
\end{gather*}
and obtains the
transition coef\/f\/icient $\gamma^A_B$ uniquely.
We have suppressed its dependence on ${\bf i}$, ${\bf j}$
in the notation.
Many remarkable properties are known for $\gamma^A_B$
including the fact $\gamma^A_B \in \Z[q]$.
See \cite[Proposition~2.3]{L} for example.

In this paper we show that the transition coef\/f\/icients
$\gamma=(\gamma^A_B)$ coincide with
the matrix elements of the intertwiner between the irreducible $A_q(\geh)$-modules labeled by two dif\/ferent reduced expressions of the longest element of the Weyl group of $\geh$.
Here $A_q(\geh)$ denotes the quantized algebra of functions  associated with
$\geh$.
It is a Hopf subalgebra of the dual $U_q(\geh)^\ast$
which has been studied from a variety of aspects.
See \cite{D86,Kas93,NYM,RTF,So1,So2,VS} for example.
Let us brief\/ly recall the most relevant result to the present paper
due to Vaksman and Soibelman \cite{So1,So2,VS}.
To each reduced expression of a (not necessarily longest) element
$w = s_{i_1}\cdots s_{i_r} \in W$,
one can associate an irreducible representation
$\pi_{\bf i}$
labeled by ${\bf i}=(i_1,\ldots, i_r)$ having the form
\begin{gather*}
\pi_{\bf i}= \pi_{i_1}\otimes \cdots \otimes \pi_{i_r}: \
A_q(\geh) \rightarrow
\mathrm{End}\big(\mathcal{F}_{q_{i_1}} \otimes \cdots \otimes
\mathcal{F}_{q_{i_r}}\big),
\end{gather*}
where each component
$\pi_i: A_q(\geh) \rightarrow \mathrm{End}(\mathcal{F}_{q_i})$ is the
{\em fundamental representation} of $A_q(\geh)$ on
the $q$-oscillator Fock space
$\mathcal{F}_{q_i} = \bigoplus_{m\ge 0}\C(q)|m\rangle$.
See Section~\ref{a:ss:gr}.
The two irreducible representations $\pi_{\bf i}$ and
$\pi_{\bf j}$ with ${\bf j}=(j_1,\ldots, j_r)$
are isomorphic if
$s_{i_1}\cdots s_{i_r}=s_{j_1}\cdots s_{j_r} \in W$
are reduced expressions (Theorem \ref{a:th:so}).
Thus one has the intertwiner
$\Phi = \Phi_{{\bf i}, {\bf j}}:
\mathcal{F}_{q_{i_1}}\otimes \cdots \otimes \mathcal{F}_{q_{i_r}}
\rightarrow
\mathcal{F}_{q_{j_1}}\otimes \cdots \otimes \mathcal{F}_{q_{j_r}}$
characterized by
\begin{gather*}
\pi_{\bf j}(g)\circ \Phi = \Phi \circ \pi_{\bf i}(g)\qquad
\forall \, g \in A_q(\mathfrak{g})
\end{gather*}
up to an overall constant.
Writing the basis of the Fock space
$\mathcal{F}_{q_{i_1}} \otimes \cdots \otimes
\mathcal{F}_{q_{i_r}}$ as
$|A\rangle = |a_1\rangle \otimes \cdots \otimes |a_r\rangle$
with $A=(a_1, \ldots, a_r) \in (\Z_{\ge 0})^r$,
we def\/ine the matrix elements of $\Phi=(\Phi^A_B)$ by
$\Phi |B\rangle = \sum_A \Phi^A_B |A\rangle$ and the
normalization $\Phi^{0,\ldots, 0}_{0,\ldots, 0}=1$.
Our main result (Theorem \ref{theorem:gamma=R})
is concerned with the longest element case $r=l$
and is stated for each pair $({\bf i}, {\bf j})$ as $\gamma^A_B = \Phi^A_B$, i.e.,
\begin{gather}\label{a:main}
\gamma= \Phi.
\end{gather}
For a convenience we also introduce the ``checked'' intertwiner
$\Phi^{\vee} = \Phi \circ \sigma$, where
$\sigma  (|a_1\rangle \otimes \cdots \otimes |a_l\rangle) =
 |a_l\rangle \otimes \cdots \otimes |a_1 \rangle$
is the reversal of the components.
